% Extracción quirúrgica del propietario V2 05b: líneas 100--359.
% El resto permanece en este archivo como procedencia, pero no se publica.
\iffalse
% Macroestructura solenoidal estrictamente interna del continuo discreto.
% Este módulo no introduce ninguna ecuación clásica ni un problema exterior.

\chapter{Internal macrostructure of balance, memory, and telescopia}
\label{chap:pdfv2-sol-interna}

\begin{statusbox}
The only input of this chapter is the discrete structure of the continuum.
The branch \(\mathrm{sol}\) is constructed before any exterior scalar evaluation
and before any classical differential formulation.  Its complete chain is
\[
 \mathcal C_{\rm cont}^{\rm disc}
 \supset\mathcal C_{\rm cont}^{\rm sol}
 \xrightarrow{\operatorname{Res}_{\rm sol}}
 \mathcal G_{\rm sol}
 \xrightarrow{\operatorname{pr}_{X,{\rm sol}}}
 X_{\chi,{\rm sol}}
 \xrightarrow{\mathcal A_{\rm sol}}
 \mathfrak M_{\rm sol}
 \xrightarrow{\mathcal P_{\rm sol}}
 \mathcal C_{\rm sol}^{\rm HMT}.
\]
Each arrow preserves as coordinates the history and the apparatus that
produces it.  No image in this chain is a quotient that forgets its source.
\end{statusbox}

\section{Proper domain prior to any subsequent evaluation}
\label{sec:pdfv2-sol-dominio}

Let
\[
 x=(x_0\xrightarrow{e_0}x_1\xrightarrow{e_1}\cdots)
 \in\mathcal C_{\rm cont}^{\rm disc}
\]
a surviving history.  Each edge \(e_j\) preserves its two sheets,
orientation, residue, quotient, and integer carry.  Let
\(N_j^+(e_j)\) and \(N_j^-(e_j)\) be the multiplicities of the two sheets, and let
\(\Delta_{\gamma_j}\in\mathbb Z\) be the oriented coefficient of the route, with the
carry of the composition already incorporated.  The internal signed increment is
\begin{equation}\label{eq:pdfv2-sol-b}
 b_j=b(e_j):=
 \Delta_{\gamma_j}\bigl(N_j^+(e_j)-N_j^-(e_j)\bigr)\in\mathbb Z.
\end{equation}
It is not obtained by evaluating an external path: it is a coordinate of the ledger of the edge itself.

Its positive and negative parts are constructed by
\begin{equation}\label{eq:pdfv2-sol-bpm}
 b_j^+:=\frac{|b_j|+b_j}{2},
 \qquad
 b_j^-:=\frac{|b_j|-b_j}{2}.
\end{equation}
Therefore,
\begin{equation}\label{eq:pdfv2-sol-bpm-identidades}
 b_j=b_j^+-b_j^-,
 \qquad
 |b_j|=b_j^++b_j^-,
 \qquad
 b_j^+b_j^-=0.
\end{equation}

The two original memoirs and their two readers are
\begin{align}
 M_0^+=M_0^-&=0,
 &M_j^+&=\sum_{0\leq i<j}b_i^+,
 &M_j^-&=\sum_{0\leq i<j}b_i^-;
 \label{eq:pdfv2-sol-memorias}\\
 D_j^0&:=M_j^+-M_j^-,
 &H_j^0&:=M_j^++M_j^-.
 \label{eq:pdfv2-sol-dh}
\end{align}
Consequently,
\begin{equation}\label{eq:pdfv2-sol-dh-diferencias}
 D_{j+1}^0-D_j^0=b_j,
 \qquad
 H_{j+1}^0-H_j^0=|b_j|.
\end{equation}
Signed memory \(D^0\) and total memory \(H^0\) remain distinct
objects; retaining only their final difference would destroy the genealogy.

The substructure \(\mathcal C_{\rm cont}^{\rm sol}\) consists of the histories
for which, in every prefix and every nonadic refinement:
\begin{enumerate}
 \item the increments \eqref{eq:pdfv2-sol-b} and the two memories
       \eqref{eq:pdfv2-sol-memorias} exist;
 \item the fine cylinders are grouped into fibres of nine extensions with the
       same projective probability law;
 \item the depth and resolution projections defined below are
       compatible with the history extensions;
 \item each loss row and each terminal column is finite at a finite
       horizon;
 \item every prefix admits at least one compatible extension to any
       subsequent horizon.
\end{enumerate}
These are conditions of membership in the branch, not silent premises
attributed to an arbitrary history.

\fi
\section{Common nonadic refinement: inclusion, balance expectation and memory}
\label{sec:pdfv2-sol-e9j9}

Let \(Y_j(x)\) be the finite set of cylinders compatible with the prefix
\(x_{\leq j}\), endowed with the projective measure \(\mu_j\).  The truncation
\[
 \tau_{j+1,j}:Y_{j+1}(x)\longrightarrow Y_j(x)
\]
has nine admissible extensions over each cylinder. Finite spaces are considered.
\[
 \mathscr H_j(x):=\ell^2\bigl(Y_j(x),\mu_j\bigr).
\]
The inclusion and the nonadic expectation are
\begin{align}
 J_{9,j}:\mathscr H_j&\longrightarrow\mathscr H_{j+1},
 &(J_{9,j}f)(y)&=f(\tau_{j+1,j}y),
 \label{eq:pdfv2-sol-j9}\\
 E_{9,j}:\mathscr H_{j+1}&\longrightarrow\mathscr H_j,
 &(E_{9,j}g)(z)&=\frac1{9}
   \sum_{\tau_{j+1,j}y=z}g(y).
 \label{eq:pdfv2-sol-e9}
\end{align}
The projective compatibility of \(\mu_{j+1}\) and \(\mu_j\) gives
\begin{equation}\label{eq:pdfv2-sol-ej-projector}
 E_{9,j}J_{9,j}=I_{\mathscr H_j},
 \qquad
 J_{9,j}^*=E_{9,j},
 \qquad
 P_{9,j}:=J_{9,j}E_{9,j}=P_{9,j}^*=P_{9,j}^2.
\end{equation}
The complementary microscopic projector is
\begin{equation}\label{eq:pdfv2-sol-qmicro-preliminar}
 Q_{9,j}:=I-P_{9,j}.
\end{equation}

If \(b_f\) is the increment on the fine fiber and
\begin{equation}\label{eq:pdfv2-sol-bcoarse}
 b_c:=E_{9,j}b_f,
\end{equation}
the elementary convexity of \(|\cdot|\) constructs, rather than postulates, the carry of
cancellation
\begin{equation}\label{eq:pdfv2-sol-kappa}
 \kappa_j
 :=\frac{E_{9,j}|b_f|-|b_c|}{2}\geq0.
\end{equation}
Indeed, using \(t^\pm=(|t|\pm t)/2\), one obtains
\begin{align}
 E_{9,j}b_f^+
 &=\frac{E_{9,j}|b_f|+E_{9,j}b_f}{2}
   =b_c^++\kappa_j,
 \label{eq:pdfv2-sol-jensen-plus}\\
 E_{9,j}b_f^-
 &=\frac{E_{9,j}|b_f|-E_{9,j}b_f}{2}
   =b_c^-+\kappa_j.
 \label{eq:pdfv2-sol-jensen-minus}
\end{align}
Both sheets lose the same amount when compressed; therefore, the signed
difference is preserved while total memory records the cancellation:
\begin{equation}\label{eq:pdfv2-sol-jensen-dh}
 E_{9,j}(b_f^+-b_f^-)=b_c,
 \qquad
 E_{9,j}(b_f^++b_f^-)=|b_c|+2\kappa_j.
\end{equation}

\section{Orthogonal factorization of the degrees of depth and of the resolution layers in \(\mathscr H_j^{\mathrm{depth}}\widehat\otimes\mathscr H_j^{\mathrm{res}}\)}
\label{sec:pdfv2-sol-micro-shell}

To prevent the same loss from being counted both as depth and as
resolution, the space at each level is typed as
\begin{equation}\label{eq:pdfv2-sol-factor-hilbert}
 \mathscr H_j
 =\mathscr H_j^{\rm depth}\widehat\otimes
  \mathscr H_j^{\rm res}.
\end{equation}
For \(j\geq1\), the first factor carries
\[
 P_j^{\rm depth}:=J_{9,j-1}E_{9,j-1}
 \quad\text{on }\mathscr H_j^{\rm depth},
\]
and \(P_0^{\rm depth}=I\) is fixed. The second factor carries the finite
shell partition inherited from the cylinder boundary. If
\(\sigma_j:Y_j^{\rm res}\to S_j\) is that partition, its operators are
\begin{align}
 J_j^{\rm res}f&=f\circ\sigma_j,
 &E_j^{\rm res}g(s)&=
 \frac1{\#\sigma_j^{-1}(s)}
 \sum_{\sigma_j(y)=s}g(y),
 \label{eq:pdfv2-sol-shell-je}\\
 P_j^{\rm res}&:=J_j^{\rm res}E_j^{\rm res}.
 \label{eq:pdfv2-sol-shell-p}
\end{align}

The three mutually orthogonal projectors are defined.
\begin{align}
 P_j^{\rm term}
 &:=P_j^{\rm depth}\otimes P_j^{\rm res},
 \label{eq:pdfv2-sol-pterm}\\
 Q_j^{\rm micro}
 &:=(I-P_j^{\rm depth})\otimes I,
 \label{eq:pdfv2-sol-qmicro}\\
 Q_j^{\rm shell}
 &:=P_j^{\rm depth}\otimes(I-P_j^{\rm res}).
 \label{eq:pdfv2-sol-qshell}
\end{align}
Then
\begin{equation}\label{eq:pdfv2-sol-ortogonalidad}
 Q_j^{\rm micro}Q_j^{\rm shell}=0,
 \qquad
 I=P_j^{\rm term}+Q_j^{\rm micro}+Q_j^{\rm shell},
\end{equation}
and, for all \(\xi\in\mathscr H_j\),
\begin{equation}\label{eq:pdfv2-sol-pitagoras}
 \|\xi\|^2
 =\|P_j^{\rm term}\xi\|^2
  +\|Q_j^{\rm micro}\xi\|^2
  +\|Q_j^{\rm shell}\xi\|^2.
\end{equation}
The term
\((I-P_j^{\rm depth})\otimes(I-P_j^{\rm res})\) does not appear a second time: it is already
contained exactly once in \(Q_j^{\rm micro}\).

\section{Cylinder extensions and full row of losses}
\label{sec:pdfv2-sol-gamma-loss}

An extension of cylinders from level \(j\) to level \(j+1\) induces the map
\begin{equation}\label{eq:pdfv2-sol-gamma}
 \Gamma_j:\mathscr H_j\longrightarrow\mathscr H_{j+1},
 \qquad
 \Gamma_j f=f\circ\tau_{j+1,j},
\end{equation}
with all leaf, orientation, carry, memory, and path labels
transported in the cylinder basis.  For \(n>j\), we write
\begin{equation}\label{eq:pdfv2-sol-gamma-composition}
 \Gamma_{j:n}:=\Gamma_{n-1}\cdots\Gamma_j,
 \qquad
 \Gamma_{j:j}:=I.
\end{equation}

On \(\mathscr H_j\), the nonnegative increments define multiplication
operators
\begin{equation}\label{eq:pdfv2-sol-loss-multipliers}
 B_j^+:=M_{\sqrt{b_j^+}},
 \qquad
 B_j^-:=M_{\sqrt{b_j^-}},
 \qquad
 K_j:=M_{\sqrt{2\kappa_j}}.
\end{equation}
Each root is taken pointwise over the finite set of cylinders; if a
coordinate is zero, its operator is zero.  The space of losses is the external
orthogonal sum
\[
 \mathscr E_j
 :=\mathscr H_j^{(+)}\oplus\mathscr H_j^{(-)}
   \oplus\mathscr H_j^{(\kappa)}
   \oplus\mathscr H_j^{({\rm micro})}
   \oplus\mathscr H_j^{({\rm shell})},
\]
and the complete row is
\begin{equation}\label{eq:pdfv2-sol-loss-row}
 L_j:\mathscr H_j\longrightarrow\mathscr E_j,
 \qquad
 L_j\xi
 =\bigl(
 B_j^+\xi,B_j^-\xi,K_j\xi,
 Q_j^{\rm micro}\xi,Q_j^{\rm shell}\xi
 \bigr).
\end{equation}
Therefore,
\begin{equation}\label{eq:pdfv2-sol-loss-norm}
 \|L_j\xi\|^2
 =\|B_j^+\xi\|^2+\|B_j^-\xi\|^2+\|K_j\xi\|^2
  +\|Q_j^{\rm micro}\xi\|^2
  +\|Q_j^{\rm shell}\xi\|^2.
\end{equation}
The positive part, negative part, cancellation carry, microstructure, and shell
each appear once. They are neither merged into an anonymous constant nor
repeated in two summands.

\section{Complete history terminal}
\label{sec:pdfv2-sol-terminal}

Let \(\operatorname{Hist}_{J}^{\rm sol}\) be the set of histories
surviving to the horizon \(J\), each with its package of thirteen
coordinates.  The terminal space is defined as the free space on these
complete records:
\begin{equation}\label{eq:pdfv2-sol-terminal-space}
 \mathscr T_J^{\rm sol}
 :=\ell^2\!\left(
 \left\{(h,D_{{\rm sol},J}(h)):
 h\in\operatorname{Hist}_{J}^{\rm sol}\right\}
 \right).
\end{equation}
The terminal evaluation
\begin{equation}\label{eq:pdfv2-sol-terminal-evaluation}
 E_J^{\rm term}:\mathscr H_J\longrightarrow\mathscr T_J^{\rm sol}
\end{equation}
sends each basic cylinder vector to the complete record of the history that
generates it and extends linearly.  In particular, the terminal retains
fiber, relations, numbers, orientation, carry, memory, boundary, route,
multisection, survival, and reader.  It is not replaced by the last visible
value of the route.

\section{Observation column, Gram operator, and telescoping}
\label{sec:pdfv2-sol-gram}

For \(0\leq j\leq J\), the complete future column is
\begin{equation}\label{eq:pdfv2-sol-g-column}
 \begin{aligned}
 G_{j;J}:\mathscr H_j&\longrightarrow
 \mathscr T_J^{\rm sol}\oplus
 \bigoplus_{n=j}^{J-1}\mathscr E_n,\\
 G_{j;J}\xi
 &:={}
 E_J^{\rm term}\Gamma_{j:J}\xi
 \ \oplus\ 
 \bigoplus_{n=j}^{J-1}L_n\Gamma_{j:n}\xi.
 \end{aligned}
\end{equation}
For \(j=J\), the sum of perdidas is empty and
\(G_{J;J}=E_J^{\rm term}\).  The Gram terminal is
\begin{equation}\label{eq:pdfv2-sol-u-gram}
 U_{j;J}:=G_{j;J}^*G_{j;J}\succeq0.
\end{equation}

Rearranging the orthogonal sum of the codomain yields the literal identity.
\begin{equation}\label{eq:pdfv2-sol-g-recursion}
 G_{j;J}\xi
 \simeq
 \bigl(L_j\xi,\,G_{j+1;J}\Gamma_j\xi\bigr).
\end{equation}
Consequently,
\begin{equation}\label{eq:pdfv2-sol-stein}
 \boxed{U_{j;J}
 =L_j^*L_j+\Gamma_j^*U_{j+1;J}\Gamma_j}
\end{equation}
and
\begin{equation}\label{eq:pdfv2-sol-stein-difference}
 U_{j;J}-\Gamma_j^*U_{j+1;J}\Gamma_j
 =L_j^*L_j\succeq0.
\end{equation}
There is no sign mutation: the loss term is on the positive side
of the identity because it arises from a sum of squares.

Iterating \eqref{eq:pdfv2-sol-stein} yields the exact telescoping identity
\begin{equation}\label{eq:pdfv2-sol-telescopia-operatorial}
 U_{j;J}
 =\Gamma_{j:J}^*(E_J^{\rm term})^*E_J^{\rm term}\Gamma_{j:J}
  +\sum_{n=j}^{J-1}
   \Gamma_{j:n}^*L_n^*L_n\Gamma_{j:n},
\end{equation}
o, in quadratic form,
\begin{equation}\label{eq:pdfv2-sol-telescopia-normas}
 \langle U_{j;J}\xi,\xi\rangle
 =\|E_J^{\rm term}\Gamma_{j:J}\xi\|^2
  +\sum_{n=j}^{J-1}\|L_n\Gamma_{j:n}\xi\|^2.
\end{equation}
The first term preserves the complete terminal future; the sum preserves,
level by level, all the losses and all the carries that led to it.

\iffalse
\section{The thirteen coordinates of the constraint}
\label{sec:pdfv2-sol-trece}

For a prefix \(x_{\leq j}\), it is defined
\begin{equation}\label{eq:pdfv2-sol-d-package}
 \begin{aligned}
 D_{{\rm sol},j}(x):=\bigl(&
 \mathcal F_{{\rm sol},j},
 \operatorname{Ext}_{\Gamma,{\rm sol},j},
 \operatorname{Rel}_{{\rm sol},j},
 \mathcal N_{{\rm sol},j},
 \operatorname{or}_{{\rm sol},j},
 \operatorname{Carr}_{{\rm sol},j},
 \operatorname{Mem}_{{\rm sol},j},\\
 &\partial_{{\rm sol},j},
 \gamma_{{\rm sol},j},
 \operatorname{Msect}_{{\rm sol},j},
 \operatorname{Surv}_{{\rm sol},j},
 \operatorname{Term}_{{\rm sol},j},
 \mathcal L_{{\rm sol},j}\bigr).
 \end{aligned}
\end{equation}
Its components are as follows.
\begin{enumerate}
 \item \(\mathcal F_{{\rm sol},j}\) contains the cylinders \(Y_j\), their projective
       measure, \(\mathscr H_j^{\rm depth}\),
       \(\mathscr H_j^{\rm res}\), their tensor products, and the five
       spaces of losses.
 \item \(\operatorname{Ext}_{\Gamma,{\rm sol},j}\) contains
       \(J_{9,j},E_{9,j},\Gamma_j,\Gamma_{j:n}\), the shell inclusions, and
       all the truncation and refinement maps.
 \item \(\operatorname{Rel}_{{\rm sol},j}\) contains
       \eqref{eq:pdfv2-sol-bpm-identidades},
       \eqref{eq:pdfv2-sol-ej-projector},
       \eqref{eq:pdfv2-sol-jensen-plus}--
       \eqref{eq:pdfv2-sol-jensen-minus},
       \eqref{eq:pdfv2-sol-ortogonalidad} and
       \eqref{eq:pdfv2-sol-stein}.
 \item \(\mathcal N_{{\rm sol},j}\) conserves
       \(j,J,9,b_j,b_j^+,b_j^-,\kappa_j\), sheet multiplicities, fiber sizes,
       and shell dimensions.
 \item \(\operatorname{or}_{{\rm sol},j}\) preserves
       \(\Delta_{\gamma_j}\), the path order, and the inversion action
       \(b_j\mapsto-b_j\), which interchanges \(b_j^+\) and \(b_j^-\) without changing
       \(|b_j|\) or \(\kappa_j\).
 \item \(\operatorname{Carr}_{{\rm sol},j}\) preserves the integer TPK carry, the composition carry and the cancellation carry \(\kappa_j\) as distinct coordinates.
 \item \(\operatorname{Mem}_{{\rm sol},j}\) is
       \((M_j^+,M_j^-,D_j^0,H_j^0)\) together with the ordered ledger of all
       preceding increments.
 \item \(\partial_{{\rm sol},j}\) contains the oriented cylinder boundary,
       the increments \(b_j\), the row \(L_j\), and the Stein differences.
 \item \(\gamma_{{\rm sol},j}\) is the complete route \(e_0e_1\cdots e_{j-1}\), not merely
its endpoints.
 \item \(\operatorname{Msect}_{{\rm sol},j}\) contains all compatible
       nonadic and shell extensions; it does not choose a prolongation.
 \item \(\operatorname{Surv}_{{\rm sol},j}\) records the horizons
       \(J\geq j\) for which \(G_{j;J}\), \(U_{j;J}\), and
       compatible terminal objects exist.
 \item \(\operatorname{Term}_{{\rm sol},j}\) contains
       \(E_J^{\rm term}\Gamma_{j:J}\), the complete record of history and
       the trece coordinates in each horizon.
 \item \(\mathcal L_{{\rm sol},j}\) internally publishes
       \((b^\pm,M^\pm,D^0,H^0,\kappa,Q^{\rm micro},Q^{\rm shell},
       L,G,U)\) and its identities.
\end{enumerate}

For an extension \(u:x_j\to x_{j'}\), the coordinated transport
\(u_{{\rm sol},j',j}\) satisfies
\begin{equation}\label{eq:pdfv2-sol-naturality-d}
 u_{{\rm sol},j',j}D_{{\rm sol},j'}(x_{j'})
 =D_{{\rm sol},j}(\tau_{j',j}x_{j'}).
\end{equation}
The equality includes all five components of \(L\), the complete terminal, and
the composition of \(\Gamma\); it is not an equality of cardinalities alone.

\section{Constraint graph and dependent object}
\label{sec:pdfv2-sol-graphs}

We define
\begin{equation}\label{eq:pdfv2-sol-g-graph}
 \mathcal G_{\rm sol}:=\operatorname{Graph}(D_{\rm sol}),
 \qquad
 \operatorname{Res}_{\rm sol}(x):=(x,D_{\rm sol}x).
\end{equation}
The first projection is inverse over the graph:
\begin{equation}\label{eq:pdfv2-sol-res-inverse}
 \pi_1\operatorname{Res}_{\rm sol}=I,
 \qquad
 \operatorname{Res}_{\rm sol}\pi_1=I_{\mathcal G_{\rm sol}}.
\end{equation}

Let \(\chi_{\rm sol}(x)\) be the complete multisection formed by all the
cylinders, refinements, shells, rows of losses, and terminals compatible
with \(x\).  Then
\begin{equation}\label{eq:pdfv2-sol-xchi}
 X_{\chi,{\rm sol}}
 :=\{(\chi_{\rm sol}(x),x,D_{\rm sol}x):
 x\in\mathcal C_{\rm cont}^{\rm sol}\},
\end{equation}
and
\begin{equation}\label{eq:pdfv2-sol-prx}
 \operatorname{pr}_{X,{\rm sol}}(x,D_{\rm sol}x)
 :=(\chi_{\rm sol}(x),x,D_{\rm sol}x).
\end{equation}
Again, the projection that preserves \((x,D_{\rm sol}x)\) is inverse on
this dependent image.

\section{Internal assembler and publisher}
\label{sec:pdfv2-sol-assembler-publisher}

The ambient space \(\mathscr M_{\rm sol}^{\rm amb}\) consists of compatible
systems
\[
 \bigl(
 \{b,b^\pm,M^\pm,D^0,H^0,\kappa\},
 \{E_9,J_9,P^{\rm term},Q^{\rm micro},Q^{\rm shell}\},
 \{\Gamma,L,E^{\rm term},G,U\}
 \bigr)
\]
with all the previous identities. The raw assembler is
\begin{equation}\label{eq:pdfv2-sol-assembler}
 \begin{aligned}
 a_{\rm sol}:X_{\chi,{\rm sol}}&\longrightarrow
 \mathscr M_{\rm sol}^{\rm amb},\\
 a_{\rm sol}(\chi(x),x,Dx)&:=
 \bigl(
 \{b,b^\pm,M^\pm,D^0,H^0,\kappa\}_x,
 \{E_9,J_9,P,Q\}_x,
 \{\Gamma,L,E^{\rm term},G,U\}_x
 \bigr).
 \end{aligned}
\end{equation}
The macrostructure preserves its antecedent through
\begin{equation}\label{eq:pdfv2-sol-macro-graph}
 \mathfrak M_{\rm sol}:=\operatorname{Graph}(a_{\rm sol}),
 \qquad
 \mathcal A_{\rm sol}(z):=(z,a_{\rm sol}z).
\end{equation}

Let \(\mathscr Y_{\rm sol}^{\rm int}\) be the type of certificates formed by
the identities
\begin{equation}\label{eq:pdfv2-sol-certificate-list}
 \begin{gathered}
 b=b^+-b^-,\qquad |b|=b^++b^-,\\
 E_9J_9=I,\qquad J_9E_9=P_9,\\
 E_9b_f^\pm=b_c^\pm+\kappa,\qquad\kappa\geq0,\\
 I=P^{\rm term}+Q^{\rm micro}+Q^{\rm shell},
 \qquad Q^{\rm micro}Q^{\rm shell}=0,\\
 U_{j;J}=L_j^*L_j+\Gamma_j^*U_{j+1;J}\Gamma_j,\\
 \langle U_{j;J}\xi,\xi\rangle
 =\|E_J^{\rm term}\Gamma_{j:J}\xi\|^2
  +\sum_{n=j}^{J-1}\|L_n\Gamma_{j:n}\xi\|^2,
 \end{gathered}
\end{equation}
together with the naturality of all arrows. The raw publisher
\begin{equation}\label{eq:pdfv2-sol-publisher}
 p_{\rm sol}^{\rm raw}:\mathfrak M_{\rm sol}
 \longrightarrow\mathscr Y_{\rm sol}^{\rm int}
\end{equation}
publish that certificate without removing the macroobject. Finally,
\begin{equation}\label{eq:pdfv2-sol-output-graph}
 \mathcal C_{\rm sol}^{\rm HMT}
 :=\operatorname{Graph}(p_{\rm sol}^{\rm raw}),
 \qquad
 \mathcal P_{\rm sol}(m):=(m,p_{\rm sol}^{\rm raw}m).
\end{equation}

\section{Solenoidal Generation Internal Theorem}
\label{sec:pdfv2-sol-theorem}

\begin{theorem}[Internal generation of the balance macrostructure]
\label{thm:pdfv2-sol-internal-generation}
For every \(x\in\mathcal C_{\rm cont}^{\rm sol}\), the chain
\[
 \mathcal C_{\rm cont}^{\rm sol}
 \xrightarrow{\operatorname{Res}_{\rm sol}}
 \mathcal G_{\rm sol}
 \xrightarrow{\operatorname{pr}_{X,{\rm sol}}}
 X_{\chi,{\rm sol}}
 \xrightarrow{\mathcal A_{\rm sol}}
 \mathfrak M_{\rm sol}
 \xrightarrow{\mathcal P_{\rm sol}}
 \mathcal C_{\rm sol}^{\rm HMT}
\]
naturally constructs a macrostructure containing both memories, the cancellation carry, the orthogonal micro/shell separation, the complete row of losses, the future column, the full-history terminal object, the Gram matrix, and the telescoping structure. The thirteen coordinates of the continuum can be recovered by successive projections from any point in the chain.
\end{theorem}

\begin{proof}
The identities \eqref{eq:pdfv2-sol-bpm-identidades} and
\eqref{eq:pdfv2-sol-dh-diferencias} are obtained directly from the
positive and negative decomposition of the integer increment; therefore the
memories are not added data.  The formulas
\eqref{eq:pdfv2-sol-j9}--\eqref{eq:pdfv2-sol-e9}, together with the projective
measure, prove \(E_9J_9=I\), \(J_9^*=E_9\), and that \(P_9\) is an
orthogonal projector.

The triangle inequality gives
\(E_9|b_f|\geq|E_9b_f|=|b_c|\), so that
\(\kappa\geq0\).  Substituting
\(b^\pm=(|b|\pm b)/2\) proves
\eqref{eq:pdfv2-sol-jensen-plus} and
\eqref{eq:pdfv2-sol-jensen-minus} without an additional hypothesis.

The tensorial definitions
\eqref{eq:pdfv2-sol-pterm}--\eqref{eq:pdfv2-sol-qshell} prove that
\(P^{\rm term}\), \(Q^{\rm micro}\), and \(Q^{\rm shell}\) are orthogonal and
sum to the identity. Consequently, row \(L_j\) counts each sector only
once.

The definition \eqref{eq:pdfv2-sol-g-column} orthogonally separates the terminal
and the losses of each level.  Separating the first summand of losses produces
\eqref{eq:pdfv2-sol-g-recursion}.  Taking the adjoint and composing proves the
Stein identity \eqref{eq:pdfv2-sol-stein}; iterating it proves
\eqref{eq:pdfv2-sol-telescopia-operatorial} and
\eqref{eq:pdfv2-sol-telescopia-normas}.  Every term is a squared norm,
so the opposite sign cannot appear.

Compatibility of cylinder truncation, expectation, projections, and
extensions \(\Gamma\) yields the naturality
\eqref{eq:pdfv2-sol-naturality-d}. Finally,
\(\mathcal G_{\rm sol}\), \(\mathfrak M_{\rm sol}\), and
\(\mathcal C_{\rm sol}^{\rm HMT}\) are graphs, and
\(X_{\chi,{\rm sol}}\) is a dependent object that explicitly preserves
its antecedent. The first projection at each stage recovers the preceding
stage. Thus no composition erases history, memory, multisection,
terminal, or reader.
\end{proof}

\section{Probative origin and forgers}
\label{sec:pdfv2-sol-provenance-falsifiers}

\paragraph{Result recovered.}
The signed increment with carry, the separation \(b=b^+-b^-\), the memories
\(M^\pm,D^0,H^0\), the pair \(E_9,J_9\), the Jensen carry \(\kappa\), cylinder
extension, the loss row, the terminal, the Gram matrix, Stein, and telescoping
belong to the already constructed architecture.

\paragraph{New Formalization.}
Explicit formalizations of that architecture are the typing prior to every
external evaluation, the factorization
\(\mathscr H^{\rm depth}\widehat\otimes\mathscr H^{\rm res}\), the orthogonal
separation preventing double counting, the free terminal on histories with
thirteen coordinates, and the four nonablative graph packagings.

\begin{criterion}[Internal falsifiers of the branch \(\mathrm{sol}\)]
\label{crit:pdfv2-sol-falsifiers}
The construction fails if any of the following occurs:
\begin{enumerate}
 \item \(b\neq b^+-b^-\) o \(|b|\neq b^++b^-\);
 \item \(E_9J_9\neq I\), the measure is not projective or \(P_9\) is not a
       projector;
 \item \(\kappa<0\) or fail
       \(E_9b_f^\pm=b_c^\pm+\kappa\);
 \item \(Q^{\rm micro}Q^{\rm shell}\neq0\), or the same sector appears in
       both summands of row \(L\);
 \item one of \(b^+,b^-,\kappa,Q^{\rm micro},Q^{\rm shell}\) is omitted or
       counted twice;
 \item the terminal preserves only the visible value and not the
history with its thirteen coordinates;
 \item the Stein identity fails or the telescoping contains a sign
       contrary to a sum of squares;
 \item an extension does not commute with \(E_9,J_9,L,G,U\);
 \item the multisection of continuations is replaced with a retrospectively chosen prolongation;
 \item any of the thirteen coordinates is omitted.
\end{enumerate}
In the last case, the binding diagnosis is: \emph{object replaced;
conclusion not transferable}.
\end{criterion}

\begin{warningbox}
This chapter concludes only the internal construction of
\(\mathcal C_{\rm sol}^{\rm HMT}\). It contains no classical differential
equation, classical initial data, map from external data, or
conclusion concerning an external problem. Those intertwiners belong to another
document and do not govern the existence of this macrostructure.
\end{warningbox}
\fi
